\chapter{Introducere}

\section{Motivație}

Când un elev lucrează unu la unu cu un profesor potrivit lui, învață mult mai bine
decât la clasă. Problema e cum îl găsește. De obicei totul se rezolvă prin
cunoștințe: întrebi pe cineva, care întreabă pe altcineva, și așa mai departe. Iar
un profesor care vrea să dea meditații n-are unde să își facă un profil, să fie
găsit de elevi și să vorbească apoi cu ei într-un singur loc.

Platformele generaliste de comunicare (rețele sociale, aplicații de mesagerie) nu oferă
structura necesară unei relații profesor-elev: nu există un mod clar de a exprima ce
predai sau ce vrei să înveți, nu există un sistem de conectare bazat pe cerere și acceptare,
iar comunicarea și partajarea materialelor didactice se realizează ad-hoc, fără un istoric
organizat pe fiecare colaborare. Există, așadar, o nevoie clară pentru o platformă dedicată,
care să structureze întregul parcurs: de la crearea unui profil, la găsirea partenerului
potrivit, la comunicarea propriu-zisă.

\section{Context și preliminarii}

Platformele educaționale existente se împart, în linii mari, în două categorii:
platforme de cursuri înregistrate, fără interacțiune personalizată, și platforme de
mediere între tutori și elevi, orientate comercial. Foarte puține pun accentul pe
comunicarea directă, în timp real, odată ce colaborarea a fost stabilită, sau pe
posibilitatea profesorului de a trimite un material de lucru chiar în conversația
în care discută cu elevul.

CourseHub este gândit ca o soluție dedicată acestei nevoi: nu un magazin de cursuri,
ci un spațiu în care oricine se poate înregistra ca persoană care vrea să predea sau
să învețe, poate căuta un partener după criterii clare (rol, nivel, materie), poate
iniția o colaborare printr-un mecanism explicit de cerere și acceptare, iar apoi poate
comunica prin chat în timp real, inclusiv prin partajarea de fișiere.

\section{Obiectivele lucrării}

Obiectivul principal al acestei lucrări este proiectarea și implementarea platformei
CourseHub, folosind ASP.NET Core 8 și tehnologii moderne asociate ecosistemului
.NET. Obiectivele specifice sunt:

\begin{itemize}
    \item construirea unui sistem de înregistrare securizat, cu verificare a adresei de
    email printr-un cod de unică folosință (OTP) și protecție a formularului de
    înregistrare printr-un mecanism de tip \emph{CAPTCHA};
    \item implementarea unui flux de completare a profilului (\emph{wizard}) prin care
    utilizatorul își alege rolul (profesor sau elev), detaliile relevante rolului ales
    (nivel de experiență, respectiv nivel de studii) și materiile de interes;
    \item dezvoltarea unui sistem de căutare și filtrare a utilizatorilor după rol, nivel
    și materie;
    \item construirea unui mecanism explicit de conectare, bazat pe cereri de colaborare
    trimise de un elev unui profesor și acceptate sau respinse de acesta;
    \item implementarea unui sistem de mesagerie în timp real, folosind SignalR,
    cu suport pentru partajarea de fișiere de către profesori către elevii cu care
    colaborează;
    \item construirea unui sistem de roluri real, bazat pe \emph{ASP.NET Identity Roles},
    și a unui panou de administrare care permite gestionarea utilizatorilor, a materiilor
    disponibile pe platformă, precum și moderarea conținutului trimis prin chat;
    \item realizarea unei interfețe complet localizate, disponibilă atât în limba română,
    cât și în limba engleză.
\end{itemize}

\section{Contribuția personală}

Toate componentele platformei CourseHub -- modelul de date, logica din controllere,
interfața, autentificarea și verificarea prin OTP, sistemul de roluri, panoul de
administrare, mesageria în timp real și localizarea -- au fost proiectate și
implementate integral ca parte a acestei lucrări, cu accent pe corectitudinea
fluxurilor de autorizare, pe validarea riguroasă a datelor introduse de utilizator
și pe o experiență de utilizare coerentă în ambele limbi.

\section{Structura lucrării}

Capitolul 2 analizează cerințele funcționale și nefuncționale, împreună cu actorii
implicați, iar capitolul 3 prezintă tehnologiile folosite. Capitolul 4 descrie
modelarea bazei de date și arhitectura sistemului, capitolul 5 detaliază
implementarea, capitolul 6 aspectele de securitate, iar capitolul 7 metodologia de
testare. Capitolul 8 conține concluziile, limitările și direcțiile de dezvoltare
ulterioară.
